\documentclass[12pt,a4paper]{article}
\usepackage[utf8]{inputenc}
\usepackage[french]{babel}
\usepackage[T1]{fontenc}
\usepackage{mathtools}
\usepackage{amsfonts}
\usepackage{amssymb}
\usepackage[pdftex]{graphicx}
\usepackage{comment}
\usepackage{siunitx}
\usepackage[colorinlistoftodos,prependcaption,textsize=tiny]{todonotes}
\usepackage[left=2cm,right=2cm,top=2cm,bottom=2cm]{geometry}
\author{Angelo Ortiz}
\title{Rapport du Projet Sorting Robot}

\begin{document}
\begin{titlepage}
  \centering
  \includegraphics[width=0.30\textwidth]{logo.jpg}\par\vspace{1cm}
  {\scshape\LARGE Sorbonne Universit\'e\par}
  \vspace{1cm}
  {\scshape\Large 2I006 : Algorithmique appliqu\'ee et structures de 
donn\'ees\par}
  \vspace{1.5cm}
  {\Large \bfseries Projet :\par}
  {\huge\bfseries Sorting Robot\par}
  \vspace{2cm}
  {\Large\itshape Kamil Rzeszutko\par}
  {\Large\itshape Angelo Ortiz\par}
  \vfill
  
  % Bottom of the page
  {\large Licence d'Informatique\par}
  {\large Ann\'ee 2017/2018\par}
\end{titlepage}
 
%\newpage
\tableofcontents
  
\newpage
  
\part*{Introduction}
\section*{Pr\'esentation}
Le projet mini projet consiste en la r\'esolution du probl\`eme du robot 
trieur...
L’objectif de ce travail est de sensibiliser les étudiants au choix de la structure de données la plus adéquate dans le cadre d’un projet défini, ici, la gestion efficace d’un conteneur de données, au travers des mesures de temps de calcul et de la mémoire utilisée pour une telle mise en place.
\section*{Aper\c{c}u}
...
\section*{L'organisation du code}
Le r\'epertoire correspondant \`a ce projet comporte plusieurs 
sous-r\'epertoires que l'on expliquera en d\'etail ci-dessous.
\subsubsection*{bin}
Ce dossier contient les fichiers ex\'ecutables des test des fonctions 
impl\'ement\'ees.
\subsubsection*{data}
Ce dossier contient les r\'esultats des diff\'erents tests de performance des 
algorithmes impl\'ement\'es.
\subsubsection*{img}
Ce dossier contient les courbes de comparaison obtenues \`a partir des donn\'ees 
du dossier {\bfseries data}.
\subsubsection*{include}
Ce dossier contient les fichiers d'en-t\^ete r\'epartis en deux sous-dossiers 
{\bfseries lib} et {\bfseries src}, chacun contenant les fichiers associ\'es au 
fichiers source des r\'epertoires du m\^eme nom.
\subsubsection*{lib}
Ce dossier contient les fichiers source fournis 
\subsubsection*{Makefile}
Ce fichier automatise la compilation des sources
\subsubsection*{obj}
Ce dossier contient les fichiers objet correspondant aux fichiers sources et aux 
fichiers ex\'ecutables
\subsubsection*{report}
Ce dossier contient ce document et les fichiers \LaTeX{} n\'ecessaires pour sa 
bonne mise en forme 
\subsubsection*{scripts}
Ce dossier contient les commandes \`a transmettre \`a {\itshape gnuplot} pour 
obtenir des courbes \`a partir des donn\'ees brutes
\subsubsection*{solutions}
Ce dossier contient les solutions obtenues, i.e.\ les cha\^ines de caract\`eres 
d\'efinissant les mouvements du robot trieur, pour les diff\'erentes fonctions 
de r\'esolution
\subsubsection*{src}
Ce dossier contient les fichiers source des fonctions de base et des fonctions 
de r\'esolution du probl\`eme du robot trieur que l'on a impl\'ement\'e
\subsubsection*{tests}
Ce dossier contient les fichiers source des tests de fonctionnement et de performance.
\addcontentsline{toc}{part}{Introduction}
\newpage

\part{Algorithme au plus proche}
Dans cette premi\`ere partie, on vous pr\'esente l'analyse faite pour chacune 
des versions impl\'ement\'ees de l'algorithme au plus proche, \`a savoir les 
versions na\"ive, circulaire, par couleur et par AVL.

\section{Version na\"ive}
\subsection*{Plus court chemin}
Soient les deux cases $(i,j)$ et $(k,l)$ dans une grille \`a $m$ 
lignes et $n$ colonnes. 
Soit la fonction \texttt{dist}$((i,j),(k,l))=|k-i|+|l-j|$. Soit la 
propri\'et\'e 
suivante $P(r), r \geq 0$:
\begin{itemize}
\item Le chemin $VH$ qui consiste \`a se d\'eplacer de $|k-i|$ cases 
verticalement vers $(k,j)$, puis de $|l-j|$ cases horizontalement vers $(k,l)$, 
et
\item le chemin $HV$ qui consiste \`a se d\'eplacer de $|l-j|$ cases 
horizontalement vers $(i,l)$, puis de $|k-i|$ cases verticalement vers $(k,l)$,
\end{itemize}
sont des plus courts chemins, o\`u $r=$ \texttt{dist}$((i,j),(k,l)) \geq 0$.

Montrons cette propri\'et\'e par r\'ecurrence faible sur $r \geq 0$.

\subsubsection*{Preuve}
\underline{Base:} Pour $r=0$, $(k,l)=(i,j)$. On se d\'eplace de $0$ 
case verticalement et horizontalement. Ainsi, on reste dans la m\^eme case. 
Donc, la propri\'et\'e est v\'erifi\'ee pour $r=0$.

\underline{Induction:} Supposons que la propri\'et\'e soit v\'erifi\'ee pour un 
$0 \leq r \leq m+n-3$ fix\'e.
Montrons que la propri\'et\'e est aussi v\'erifi\'ee pour $r+1$.

Soient $p=|k-i|, q=|j-l| \in \mathbb{N}$, tels que $r+1=p+q$. Puisque 
$r+1 \geq1$, au moins l'une des variables est non nulle. Supposons sans perte 
de g\'en\'eralit\'e que $p > 0$, i.e.\ $p-1 \geq 0$.

Trois cas sont possibles:
\begin{enumerate}
  \item \underline{$i=0$:} 

  Tout d'abord, $p=|k-0|=k$. Puis, pour aller de la case $(0,j)$ vers la case 
$(k,l)$ avec un premier   d\'eplacement vertical d'une case, il faut passer par 
la case $(1,j)$. On a que $|k-1|=p-1$ et, en 
cons\'equence, \texttt{dist}$((1,j),(k,l))=(p-1)+q=r$. Par 
hypoth\`ese de r\'ecurrence, $P(r)$ est v\'erifi\'ee. De ces deux faits, le 
chemin se d\'epla\c{c}ant de $|k|$ cases verticalement vers $(k,j)$, puis de 
$|l-j|$ cases horizontalement vers $(k,l)$ est un plus court chemin.
  \item \underline{$i=m-1$ :} 

  Tout d'abord, $p=|k-(m-1)|=m-k-1$. Puis, pour aller de la case $(m-1,j)$ vers 
la case $(k,l)$ avec un premier d\'eplacement vertical d'une case, il faut 
passer par la case $(m-2,j)$. On a que $|k-(m-2)|=(m-2)-k=p-1$ et, en 
cons\'equence, \texttt{dist}$((m-2,j),(k,l))=(p-1)+q=r$. Par hypoth\`ese de 
  r\'ecurrence, $P(r)$ est v\'erifi\'ee. De ces deux faits, le chemin se 
  d\'epla\c{c}ant de $|k-(m-1)|$ cases verticalement vers $(k,j)$, puis de 
$|l-j|$ cases horizontalement vers $(k,l)$ est un plus court chemin.
  \item \underline{$0 < i <  m-1$ :}

  Pour aller de la case $(i,j)$ vers la case $(k,l)$ avec un premier 
d\'eplacement vertical d'une case, il faut passer par la case $(h,j)$, 
o\`u $h \in \{i-1, i+1\}$, telle que $0 \leq h < m$ et $|k-h|=p-1$. On a alors 
que \texttt{dist}$((h,j),(k,l))=(p-1)+q=r$. Par hypoth\`ese de 
r\'ecurrence, $P(r)$ est v\'erifi\'ee. Donc, en particulier, le chemin se 
d\'epla\c{c}ant de $|k-h|$ cases verticalement vers $(k,j)$, puis de $|l-j|$ 
cases horizontalement vers $(k,l)$ est un plus court chemin.
\end{enumerate}
On fait de m\^eme pour $q>0$ avec un traitement non plus sur les lignes, mais 
sur les colonnes de la grille, et on obtient que le chemin se d\'epla\c{c}ant de 
$|l-j|$ cases horizontalement vers $(i,l)$, puis de $|k-i|$ cases verticalement 
vers $(k,l)$ est un plus court chemin

\underline{Conclusion} : 
\begin{equation*}
  \left .\begin{array}{l}
	P(0) \text{ vraie } \\
	\forall r \in \{0,\dotsc, m+n-3\}, [ P(r) \implies P(r+1) 
	]
  \end{array} \right \}
  \left .\begin{array}{l}
	\forall r \in \{0,\dotsc, m+n-2\}, \\
	\text{les chemins } VH \text{ et } HV\\
	\text{sont des plus courts chemins.}
  \end{array}\right .
\end{equation*}

\subsection*{Fonctions auxiliaires}
\begin{itemize}
  \item La fonction \texttt{estCaseNoire} nous permet de savoir si une case est 
noire, i.e.\ si elle porte une pi\`ece de m\^eme couleur.
  \item La fonction \texttt{estPieceNoire} nous permet de savoir si une pi\`ece 
est non noire, i.e.\ si sa couleur est diff\'erente de $-1$.
  \item La fonction \texttt{robotPortePiece} nous permet de savoir si le robot 
porte une pi\`ece, i.e.\ si la couleur de la ``pi\`ece'' du robot est 
diff\'erente de $-1$.
  \item La fonction \texttt{couleurPieceRobot} nous permet de savoir dans le 
cas o\`u le robot porte une pi\`ece, sa couleur.
\end{itemize}

\subsection*{Analyse de la complexit\'e}
Pour l'analyse de la complexit\'e des quatre impl\'ementations, on 
commence par les morceaux de code communs aux toutes. Dans le pire cas, 
il n'y a aucune case noire dans la grille, i.e.\ il y a $O(n^2)$ pi\`eces \`a 
traiter, ce qui appara\^it sous la forme d'une boucle \texttt{while} dans le 
code fourni, en l'occurrence la boucle principale des fonctions. De 
plus, la recherche d'une pi\`ece non noire aux alentours du robot est faite au 
tout d\'ebut de l'algorithme et \`a chaque fois que l'on ferme un 
circuit : cela d\'epend fortement du nombre de cycles dans la grille. Par 
cons\'equent, ceci n'est pas pris en compte pour l'analyse de la complexit\'e et 
celle-ci se r\'eduit \`a la recherche de pi\`eces. Par ailleurs, on a un 
param\`etre $\alpha = \left \lceil \frac{nm}{c} 
\right \rceil$, o\`u $m$, $n$ et $c$ correspondent respectivement au nombre 
de lignes, de colonnes et de couleurs de la grille. Ici, $\alpha = \left \lceil 
\frac{n^2}{c} \right \rceil$.\par
Les deux premi\`eres impl\'ementations de l'algorithme au plus proche ont une 
complexit\'e en $O(n^4)$. Pour chaque pi\`ece on cherche la case la plus proche 
dont la couleur co\"incide avec la sienne. L'impl\'ementation na\"ive de 
l'algorithme fait une recherche sur toute la grille, ce qui est de l'ordre de 
$\Theta(n^2)$. Quant \`a la deuxi\`eme, la recherche est circulaire et on a donc 
une complexit\'e en $O(n^2)$.

Il est imp\'eratif de remarquer que la recherche circulaire est 
plus rapide au fur et \`a mesure que le param\`etre $\alpha$ augmente. \par
On a compar\'e les temps d'ex\'ecution des versions na\"ive et circulaire pour 
des grilles carr\'ees de taille variable : la premi\`ere version met 
\SI{31,48}{\second} pour une grille de longueur 110, tandis que la deuxi\`eme 
met \SI{32,25}{\second} pour une grille de longueur 155.

Les exp\'erimentations ont \'et\'e faites avec un nombre de couleurs \'egal \`a 
$\left \lfloor \frac{n}{2} \right \rfloor$, ce qui donne $\alpha \approx 2n$. 
Comme dit pr\'ec\'edemment, ce param\`etre est important pour la dur\'ee 
effective de la recherche circulaire. C'est pourquoi on a obtenu des meilleurs 
r\'esultats pour la deuxi\`eme version alors que la complexit\'e pire-cas est 
la m\^eme pour toutes les deux.

\section{Version circulaire}
\subsection*{Analyse de la complexit\'e}
...

\section{Version par couleur}
\subsection*{Analyse de la complexit\'e}
L'analyse de la complexit\'e de l'impl\'ementation par couleur de 
l'algorithme au plus proche se fera en deux parties. Dans un premier temps, 
on d\'etermine une borne sup\'erieure de la complexit\'e. En effet, le 
robot porte une pi\`ece \`a presque tout moment du jeu et cherche alors la 
case de m\^eme couleur la plus proche de lui. On a un acc\`es en $\Theta(1)$ \`a 
la liste cha\^in\'ee correspondant \`a la couleur de la pi\`ece et on est 
oblig\'e de parcourir toute la liste car il n'y a pas d'ordre d\'efini 
lors de l'insertion. Soit $l$ la longueur de la liste cha\^in\'ee. On a donc que 
la complexit\'e de la recherche de ladite case est en $\Theta(l)$. Or on a $l 
\leq n^2$. On en d\'eduit que le traitement d'une pi\`ece est de l'ordre de 
$O(n^2)$, ce qui donne une complexit\'e en $O(n^4)$ pour cette version. 
Nonobstant, on a n\'eglig\'e un facteur cl\'e de cette version : le param\`etre 
$\alpha$ repr\'esentant le nombre maximal de pi\`eces d'une m\^eme couleur et 
ainsi celui des cases de m\^eme couleur. \par
Dans l'\'etape suivante, on tient compte de ce param\`etre dans l'analyse 
et on remarque qu'il remplace $n^2$ dans le paragraphe ci-dessus. De ce fait, 
on obtient que la complexit\'e de la recherche dans la liste cha\^in\'ee est en 
$O(\alpha)$, ce qui donne finalement une complexit\'e en $O(\alpha n^2)$ pour la 
version par couleur de l'algorithme.

{\bfseries N.B.} L'insertion de chacune des $n^2$ cases dans la liste 
cha\^in\'ee correspondante est en $\Theta(1)$, ce qui donne une complexit\'e de 
l'insertion en $\Theta(n^2)$. On remarque ainsi que les complexit\'es donn\'ees 
ci-dessus l'emportent \`a l'\'egard de celle de l'initialisation des listes.
\begin{comment}
  Puisque la m\'emoire allou\'ee est toujours 
  la m\^eme ($mn$ cellules), on a plut\^ot int\'er\^et \`a diminuer $\alpha$, 
  i.e.\ augmenter le nombre de couleurs pr\'esentes dans la grille.
\end{comment}

\section{Version par AVL}
\subsection*{Analyse de la complexit\'e}
La derni\`ere version de cette algorithme r\'epose sur l'utilisation des 
arbres AVL. On a rang\'e les cases dans une matrice d'arbres : chaque ligne 
correspond \`a une couleur distincte et chaque colonne correspond \`a une ligne 
de la grille du jeu. Etant donn\'e que l'on travaille avec une matrice, 
l'acc\`es \`a un certain arbre est en $\Theta(1)$. Cependant, la recherche de 
la case la plus proche dans une ligne donn\'ee est en $O(h(T))$, o\`u $h(T)$ 
est la hauteur de l'arbre $T$ correspondant. On sait d'ailleurs que $h(T) \leq 
\left \lceil \log_2 n(T) \right \rceil$, o\`u $n(T)$ est le nombre de 
n\oe uds dans l'arbre $T$. Or on a $n(T) = \min(n,\alpha)$. 
Ceci nous donne une complexit\'e en $O(\log \min(n,\alpha))$. On en d\'eduit 
que la recherche sur les $n$ lignes de la grille est en $O(n\log 
\min(n,\alpha))$. On obtient finalement une complexit\'e en $O(n^3\log 
\min(n,\alpha))$ pour cette derni\`ere impl\'ementation de l'algorithme au plus 
proche.

{\bfseries N.B.} L'insertion de chacune des $n^2$ cases dans l'arbre 
AVL correspondant est en $O(\log \min(n,\alpha))$, ce qui donne une 
complexit\'e de l'insertion en $O(n^2\log \min(n,\alpha))$. On remarque ainsi 
que la complexit\'e donn\'ee ci-dessus l'emporte \`a l'\'egard de celle de 
l'initialisation des arbres. \par
Comme mentionn\'e ci-dessus, les exp\'erimentations ont \'et\'e faites avec 
$\alpha \approx 2n$, ce qui transforme les complexit\'es donn\'ees 
pr\'ec\'edemment en $O(n^3)$ pour la version par couleur et $O(n^3\log n)$ pour 
la version par AVL. C'est pourquoi on a obtenu des meilleurs r\'esultats pour 
la premi\`ere de ces deux versions.
\todo{Les r\'esultats empiriques ne collent pas \`a l'analyse de la 
complexit\'e}

\newpage

\part{Une case par couleur - Graphes}
Dans cette deuxi\`eme partie, on vous d\'etaille le fonctionnement des 
proc\'edures de r\'esolution du probl\`eme du robot trieur \`a l'aide des 
graphes.

\newpage

\part{Comparaison et performances}
Dans cette troisi\`eme partie, on vous pre\'esente les r\'esultats des 
diff\'erents test de performance r\'ealis\'es.

\newpage

\part*{Conclusion}
...
\addcontentsline{toc}{part}{Conclusion}

\end{document}
